% =============================================================================
% arXiv Submission Format
% The "For You" Page Problem: Explicit Content Exposure and Mental Health
% Concerns in Short-Form Video Platforms
% A Study from Nepal's Developing Digital Economy
% =============================================================================
% arXiv identifier: [To be assigned]
% Primary category: cs.SI (Social and Information Networks)
% Secondary: cs.CY (Computers and Society), cs.CL (Computation and Language)
% =============================================================================

\documentclass[11pt,a4paper]{article}

% =============================================================================
% PACKAGES (arXiv compatible)
% =============================================================================
\usepackage[utf8]{inputenc}
\usepackage[T1]{fontenc}
\usepackage{times}
\usepackage[margin=1in]{geometry}
\usepackage{setspace}
\usepackage{graphicx}
\usepackage{float}
\usepackage{booktabs}
\usepackage{array}
\usepackage{multirow}
\usepackage{tabularx}
\usepackage{amsmath,amssymb}
\usepackage{hyperref}
\usepackage{xcolor}
\usepackage{caption}
\usepackage{subcaption}
\usepackage{natbib}
\usepackage{enumitem}
\usepackage{url}

% =============================================================================
% CONFIGURATION
% =============================================================================
\onehalfspacing

% Hyperlink configuration
\hypersetup{
    colorlinks=true,
    linkcolor=blue,
    citecolor=blue,
    urlcolor=blue,
    pdfauthor={Anonymous},
    pdftitle={The For You Page Problem: Nepal Study}
}

% Figure path (relative for arXiv)
\graphicspath{{figures/}}

% Table column types
\newcolumntype{L}[1]{>{\raggedright\arraybackslash}p{#1}}
\newcolumntype{C}[1]{>{\centering\arraybackslash}p{#1}}
\newcolumntype{R}[1]{>{\raggedleft\arraybackslash}p{#1}}

% =============================================================================
% TITLE
% =============================================================================
\title{\textbf{The ``For You'' Page Problem: Explicit Content Exposure and Mental Health Concerns in Short-Form Video Platforms}\\[0.3cm]
\large A Secondary Data Analysis from Nepal's Developing Digital Economy}

\author{
    [Author Name]\\
    Independent Researcher\\
    Kathmandu, Nepal\\
    \texttt{[email@example.com]}
}

\date{January 2026}

% =============================================================================
% DOCUMENT
% =============================================================================
\begin{document}

\maketitle

% =============================================================================
% ABSTRACT
% =============================================================================
\begin{abstract}
\noindent
\textbf{Background:} Short-form video platforms (TikTok, Instagram Reels, YouTube Shorts) have experienced unprecedented growth globally, with particularly rapid adoption in South Asian developing economies like Nepal. However, concerns regarding algorithmic exposure to inappropriate content and mental health impacts remain understudied in these contexts, where digital literacy and regulatory frameworks are still evolving.

\textbf{Objective:} This study examines public discourse surrounding explicit content exposure, age-related concerns, and mental health implications of short-form video platforms through large-scale secondary data analysis, with particular attention to the unique challenges facing developing digital economies.

\textbf{Methods:} We collected 221,764 user-generated records from five sources: Google Play Store reviews (160,000), Apple App Store reviews (12,086), Reddit posts (48,606), Trustpilot reviews (800), and YouTube comments (18,539). After relevance filtering and deduplication, 8,435 records mentioning content concerns, age-related issues, or mental health were analyzed using VADER sentiment analysis and Latent Dirichlet Allocation (LDA) topic modeling.

\textbf{Results:} Analysis revealed highly polarized sentiment (45.0\% negative, 42.3\% positive, 12.7\% neutral) with platform-specific variations. Instagram exhibited the most negative sentiment ($M = -0.096$), while TikTok showed slightly positive sentiment ($M = 0.025$) despite receiving the most discussion (49.6\% of relevant content). Five dominant themes emerged: (1) child safety and parental concerns, (2) algorithmic content recommendation issues, (3) personal usage experiences, (4) account moderation and banning, and (5) mental health and attention span concerns.

\textbf{Conclusions:} Public discourse reveals significant concerns about algorithmic content exposure affecting minors, with platform-specific sentiment patterns suggesting differential user experiences. The findings highlight the need for improved content moderation, transparent algorithmic practices, and age-appropriate safeguards---particularly critical in developing economies like Nepal where digital adoption is outpacing regulatory and educational infrastructure.

\vspace{0.3cm}
\noindent\textbf{Keywords:} TikTok, Instagram Reels, YouTube Shorts, content moderation, algorithmic recommendation, mental health, child safety, sentiment analysis, topic modeling, Nepal, South Asia, developing economies
\end{abstract}

\newpage

% =============================================================================
% 1. INTRODUCTION
% =============================================================================
\section{Introduction}

\subsection{Background}

The proliferation of short-form video content has fundamentally transformed digital media consumption patterns over the past five years. Platforms such as TikTok, Instagram Reels, YouTube Shorts, and Facebook Reels have accumulated billions of users worldwide, with particularly high adoption rates among adolescents and young adults \citep{anderson2023teens}. These platforms employ sophisticated algorithmic recommendation systems---commonly referred to as ``For You'' pages (FYP)---that curate personalized content feeds based on user behavior, engagement patterns, and demographic signals.

While these algorithms optimize for user engagement and platform retention, growing concerns have emerged regarding their potential to expose users, particularly minors, to inappropriate, explicit, or harmful content \citep{nguyen2025feeds}. The landmark meta-analysis by \citet{nguyen2025feeds}, encompassing 98,299 participants across 71 studies, identified content-specific effects as a critical research gap in understanding social media's psychological impacts.

The rapid rise of short-form video platforms represents a paradigm shift in content consumption. TikTok alone reported over 1.5 billion monthly active users in 2024, with approximately 60\% aged between 13-24 years \citep{datareportal2024}. This demographic concentration raises particular concerns regarding content appropriateness and developmental impacts.

\subsection{The Nepal Context: A Developing Digital Economy}

This research is particularly motivated by the unique challenges facing \textbf{Nepal and similar developing digital economies} in South Asia. Nepal has witnessed explosive growth in internet and smartphone penetration over the past decade:

\begin{itemize}[itemsep=0.3em]
    \item Internet penetration increased from 17\% (2015) to over 65\% (2024)
    \item Mobile data prices among the lowest in the region, facilitating rapid adoption
    \item Youth population (15-24 years) comprises approximately 20\% of the population
    \item TikTok and Instagram have become dominant platforms, particularly among urban youth
    \item Limited digital literacy education in schools and communities
    \item Nascent regulatory framework for online content moderation
\end{itemize}

This context creates a ``perfect storm'' where:
\begin{enumerate}[itemsep=0.2em]
    \item Rapid platform adoption outpaces digital literacy development
    \item Parents often have less digital familiarity than their children
    \item Content moderation systems may not account for local languages and cultural contexts
    \item Regulatory mechanisms lag behind technological deployment
    \item Economic constraints limit access to premium safety features
\end{enumerate}

While most existing research on social media effects originates from Western, educated, industrialized, rich, and democratic (WEIRD) populations, the experiences and concerns of users in developing economies remain critically understudied \citep{pew2023global}.

\subsection{The Research Gap}

Despite extensive research on social media's general effects on mental health, several gaps persist:

\begin{enumerate}[leftmargin=*, itemsep=0.5em]
    \item \textbf{Content-Specific Analysis:} Most studies examine platform usage broadly rather than specific content exposure patterns \citep{valkenburg2022social}
    \item \textbf{Algorithmic Concerns:} Limited research addresses how recommendation algorithms specifically contribute to inappropriate content exposure \citep{ribeiro2020auditing}
    \item \textbf{Geographic Context:} Research from developing digital economies, particularly South Asia, remains virtually non-existent
    \item \textbf{Multi-Platform Comparison:} Few studies compare user concerns across competing short-form video platforms
    \item \textbf{Non-WEIRD Populations:} The vast majority of social media research originates from Western contexts
\end{enumerate}

\subsection{Research Questions}

This study addresses the following research questions:

\begin{itemize}[leftmargin=*, itemsep=0.3em]
    \item[\textbf{RQ1:}] What is the overall sentiment of public discourse regarding explicit content and safety concerns on short-form video platforms?
    \item[\textbf{RQ2:}] How do sentiment patterns differ across platforms (TikTok, Instagram, YouTube, Facebook)?
    \item[\textbf{RQ3:}] What are the dominant themes in user-generated content discussing short-form video platform concerns?
    \item[\textbf{RQ4:}] What specific concerns do users express regarding children, algorithms, and mental health impacts?
    \item[\textbf{RQ5:}] What implications do these findings have for developing digital economies like Nepal?
\end{itemize}

\subsection{Contributions}

This study contributes to the literature by:
\begin{itemize}[leftmargin=*, itemsep=0.3em]
    \item Providing large-scale analysis of 221,764 user-generated records across multiple platforms
    \item Identifying platform-specific sentiment patterns related to content concerns
    \item Extracting dominant discourse themes through unsupervised topic modeling
    \item Offering empirical evidence to inform platform policy and parental guidance
    \item \textbf{Providing one of the first studies addressing short-form video concerns with relevance to South Asian developing economies}
    \item Contributing to the limited literature on social media effects outside WEIRD populations
\end{itemize}

% =============================================================================
% 2. LITERATURE REVIEW
% =============================================================================
\section{Literature Review}

\subsection{Short-Form Video Platforms and Youth}

Short-form video platforms have become the dominant form of social media engagement among Generation Z and younger millennials. According to the Common Sense Media census, teens spend an average of 4.8 hours daily on screen media, with video content comprising the largest share \citep{rideout2022common}. The ``endless scroll'' design and algorithmic personalization create highly engaging experiences that have raised concerns about compulsive use patterns \citep{alter2017irresistible}.

In Nepal and South Asia, platform adoption has been even more dramatic due to:
\begin{itemize}[itemsep=0.2em]
    \item Affordable mobile data plans making video streaming accessible
    \item Strong youth demographic (median age in Nepal: 24 years)
    \item Cultural affinity for visual and musical content
    \item Limited alternative entertainment options in rural areas
    \item Social media as primary news and information source
\end{itemize}

\subsection{Algorithmic Content Curation}

The ``For You'' page algorithm represents a significant departure from traditional social media feeds. Unlike chronological or friend-based content delivery, FYP algorithms prioritize content predicted to maximize engagement, regardless of creator-viewer relationships \citep{tiktok2023safety}. This approach has been criticized for creating ``filter bubbles'' and potentially accelerating exposure to progressively extreme or inappropriate content \citep{pariser2011filter, ribeiro2020auditing}.

Research by \citet{bakshy2015exposure} demonstrated that algorithmic curation significantly shapes content exposure patterns, though the specific mechanisms vary across platforms. The opacity of these algorithms has drawn regulatory scrutiny and public concern \citep{gorwa2020algorithmic}.

\subsection{Mental Health Implications}

Recent research has documented associations between short-form video consumption and various psychological outcomes. Table~\ref{tab:literature_comparison} summarizes key findings from major studies in the field.

\begin{table}[H]
\centering
\caption{Comparison with Prior Research on Social Media and Mental Health}
\label{tab:literature_comparison}
\small
\begin{tabular}{L{2.5cm}L{2cm}C{1.5cm}L{3cm}L{3.5cm}}
\toprule
\textbf{Study} & \textbf{Platform(s)} & \textbf{N} & \textbf{Method} & \textbf{Key Finding} \\
\midrule
\citet{nguyen2025feeds} & Multiple & 98,299 & Meta-analysis & Small negative effect ($r = -0.08$) on well-being \\
\addlinespace
\citet{twenge2018increases} & General SM & 506,820 & Survey analysis & Increased depression after 2010, linked to SM \\
\addlinespace
\citet{orben2019association} & General SM & 355,358 & Large-scale survey & Effect size comparable to eating potatoes \\
\addlinespace
\citet{kelly2019social} & Multiple & 10,904 & Cohort study & SM use linked to poor sleep, cyberbullying \\
\addlinespace
\citet{coyne2020does} & General SM & 500 & 8-year longitudinal & No significant effect on depression/anxiety \\
\addlinespace
\citet{su2021fmri} & TikTok & 40 & fMRI study & Reward system activation during FYP viewing \\
\addlinespace
\textbf{This Study} & TikTok, IG, YT, FB & 8,435 & NLP/Sentiment & 45\% negative sentiment, platform differences \\
\bottomrule
\end{tabular}
\end{table}

Key documented concerns include:

\begin{itemize}[leftmargin=*, itemsep=0.3em]
    \item \textbf{Attention span:} Studies suggest that rapid-fire content consumption may impact sustained attention capabilities \citep{su2021fmri}
    \item \textbf{Body image:} Exposure to idealized content correlates with body dissatisfaction, particularly among adolescent females \citep{fardouly2022social, tiggemann2018media}
    \item \textbf{Addictive patterns:} The infinite scroll and variable reward mechanisms exhibit similarities to gambling addiction models \citep{montag2021addictive, andreassen2016bergen}
\end{itemize}

\subsection{Content Moderation Challenges}

Platform content moderation faces inherent tensions between user expression and harmful content prevention, algorithmic efficiency and human oversight, global standards and cultural context variability, and age verification and privacy concerns \citep{gillespie2018custodians}. The scale of content uploaded daily---TikTok alone processes millions of videos---makes comprehensive human review infeasible, leading to increased reliance on automated systems with known limitations \citep{gorwa2020algorithmic}.

\textbf{These challenges are amplified in developing economies} where:
\begin{itemize}[itemsep=0.2em]
    \item Local languages (e.g., Nepali, regional dialects) may not be well-supported by moderation AI
    \item Cultural context for ``inappropriate'' content varies significantly
    \item Limited local content moderation staff
    \item Users may lack awareness of reporting mechanisms
\end{itemize}

\subsection{Theoretical Framework}

This study draws upon multiple theoretical perspectives (Table~\ref{tab:theoretical_framework}):

\begin{table}[H]
\centering
\caption{Theoretical Framework for Understanding Short-Form Video Effects}
\label{tab:theoretical_framework}
\small
\begin{tabular}{L{3.5cm}L{5cm}L{5cm}}
\toprule
\textbf{Theory} & \textbf{Core Concept} & \textbf{Application to This Study} \\
\midrule
Social Learning Theory \citep{bandura1977social} & Behavior learned through observation and modeling & Users model attitudes/behaviors seen in explicit content \\
\addlinespace
Uses \& Gratifications \citep{katz1973uses} & Media use driven by specific needs & Different age groups seek content for different reasons \\
\addlinespace
Cultivation Theory \citep{gerbner1998cultivation} & Long-term exposure shapes worldview & Repeated explicit content normalizes inappropriate material \\
\addlinespace
Media Effects Model & Direct and indirect pathways of influence & Sentiment reflects perceived effects on users \\
\bottomrule
\end{tabular}
\end{table}

% =============================================================================
% 3. METHODOLOGY
% =============================================================================
\section{Methodology}

\subsection{Research Design}

This study employed a cross-sectional, secondary data analysis approach, examining publicly available user-generated content from multiple online sources. This methodology was selected for several reasons:

\begin{enumerate}[leftmargin=*, itemsep=0.3em]
    \item \textbf{Ecological validity:} User reviews and social media posts represent organic, unsolicited expressions of platform experiences
    \item \textbf{Scale:} Secondary data enables analysis of substantially larger samples than primary survey research
    \item \textbf{Accessibility:} Public data sources eliminate the ethical complexities of directly researching minor populations
    \item \textbf{Cost-effectiveness:} Zero-budget research approach suitable for independent researchers in resource-constrained settings
    \item \textbf{Relevance to Nepal:} While not Nepal-specific, findings provide baseline for understanding global discourse that shapes platform policies affecting Nepali users
\end{enumerate}

\subsection{Data Sources}

Data were collected from five publicly accessible sources (Table~\ref{tab:data_sources}):

\begin{table}[H]
\centering
\caption{Data Sources and Collection Methods}
\label{tab:data_sources}
\begin{tabular}{llll}
\toprule
\textbf{Source} & \textbf{Type} & \textbf{Collection Method} & \textbf{Records} \\
\midrule
Google Play Store & App reviews & google-play-scraper API & 160,000 \\
Apple App Store & App reviews & RSS feed scraping & 12,086 \\
Reddit & Forum posts & Public JSON API & 48,606 \\
Trustpilot & Platform reviews & Web scraping & 800 \\
YouTube & Video comments & youtube-comment-downloader & 18,539 \\
\midrule
\textbf{Total} & & & \textbf{221,764} \\
\bottomrule
\end{tabular}
\end{table}

\subsection{Target Applications}

Reviews and discussions were collected for four major short-form video platforms:
\begin{itemize}[leftmargin=*, itemsep=0.2em]
    \item TikTok (com.zhiliaoapp.musically) --- Most popular in Nepal among youth
    \item Instagram (com.instagram.android) --- Growing rapidly, especially urban areas
    \item YouTube (com.google.android.youtube) --- Established platform, Shorts gaining traction
    \item Facebook (com.facebook.katana) --- Reels feature, popular among older demographics
\end{itemize}

\subsection{Keyword Filtering}

To identify relevant content, records were filtered using keyword matching across five categories (Table~\ref{tab:keywords}):

\begin{table}[H]
\centering
\caption{Keyword Categories for Relevance Filtering}
\label{tab:keywords}
\small
\begin{tabular}{lL{10cm}}
\toprule
\textbf{Category} & \textbf{Keywords} \\
\midrule
Explicit Content & inappropriate, explicit, sexual, nudity, nude, pornography, porn, naked, nsfw, adult content \\
\addlinespace
Age Concerns & kid, child, children, teen, teenager, minor, young, daughter, son, parent, underage \\
\addlinespace
Mental Health & addiction, addicted, mental health, depression, anxiety, brain rot, attention span, dopamine, toxic \\
\addlinespace
Algorithm & algorithm, fyp, for you page, recommended, keeps showing, feed, suggestion \\
\addlinespace
Parental Controls & parental, control, filter, restrict, block, protect, safety, ban \\
\bottomrule
\end{tabular}
\end{table}

\subsection{Data Processing Pipeline}

The data processing pipeline included:

\begin{enumerate}[leftmargin=*, itemsep=0.3em]
    \item \textbf{Standardization:} Column harmonization across heterogeneous data sources
    \item \textbf{Deduplication:} Removal of duplicate records based on text content
    \item \textbf{Relevance filtering:} Retention of records containing at least one keyword
    \item \textbf{Text cleaning:} Removal of URLs, special characters, and excessive whitespace
    \item \textbf{Language detection:} Identification of primary language using langdetect
\end{enumerate}

\textbf{Final analytical sample:} 8,435 relevant records (3.8\% relevance rate)

\subsection{Sentiment Analysis}

Sentiment analysis was conducted using VADER (Valence Aware Dictionary and sEntiment Reasoner), a lexicon and rule-based sentiment analysis tool specifically designed for social media text \citep{hutto2014vader}. VADER produces a compound score ranging from $-1$ (most negative) to $+1$ (most positive).

Sentiment categories were assigned as follows:
\begin{itemize}[leftmargin=*, itemsep=0.2em]
    \item \textbf{Positive:} compound score $> 0.05$
    \item \textbf{Neutral:} $-0.05 \leq$ compound score $\leq 0.05$
    \item \textbf{Negative:} compound score $< -0.05$
\end{itemize}

\subsection{Topic Modeling}

Latent Dirichlet Allocation (LDA) was employed to identify latent themes within the corpus \citep{blei2003latent}. The modeling process included:

\begin{enumerate}[leftmargin=*, itemsep=0.3em]
    \item \textbf{Preprocessing:} Tokenization, stopword removal, lemmatization
    \item \textbf{Dictionary creation:} Filtering terms appearing in $<$5 or $>$50\% of documents
    \item \textbf{Model selection:} Coherence score optimization across $k = \{5, 7, 10, 12, 15\}$ topics
    \item \textbf{Model training:} 5 topics (optimal), 20 passes, using Gensim library
    \item \textbf{Coherence evaluation:} $C_v$ coherence score for model quality assessment \citep{mimno2011optimizing}
\end{enumerate}

\subsection{Ethical Considerations}

This study analyzed publicly available data only. All user identifiers were removed prior to analysis. No direct contact with human subjects occurred. The research was conducted in accordance with platform terms of service and ethical guidelines for internet research \citep{zimmer2010but, salganik2019bit}.

% =============================================================================
% 4. RESULTS
% =============================================================================
\section{Results}

\subsection{Sample Characteristics}

The final analytical sample comprised 8,435 relevant records after filtering and deduplication from an initial pool of 221,764 collected records (3.8\% relevance rate).

\subsubsection{Distribution by Platform}

TikTok dominated the discussion of content concerns, comprising nearly half of all relevant records (Table~\ref{tab:platform_distribution}, Figure~\ref{fig:platform_distribution}).

\begin{table}[H]
\centering
\caption{Distribution of Relevant Records by Platform}
\label{tab:platform_distribution}
\begin{tabular}{lrr}
\toprule
\textbf{Platform} & \textbf{n} & \textbf{\%} \\
\midrule
TikTok & 4,181 & 49.6 \\
Social Media General & 1,654 & 19.6 \\
Instagram & 1,011 & 12.0 \\
YouTube & 941 & 11.2 \\
Facebook & 648 & 7.7 \\
\midrule
\textbf{Total} & \textbf{8,435} & \textbf{100.0} \\
\bottomrule
\end{tabular}
\end{table}

\begin{figure}[H]
    \centering
    \includegraphics[width=0.75\textwidth]{platform_distribution.pdf}
    \caption{Distribution of relevant records across platforms. TikTok dominates the discourse with 49.6\% of all content-concern-related discussions, reflecting its popularity particularly among younger users in developing markets like Nepal.}
    \label{fig:platform_distribution}
\end{figure}

\subsubsection{Distribution by Data Source}

YouTube comments provided the largest share of relevant records, followed by Google Play Store reviews (Table~\ref{tab:source_distribution}).

\begin{table}[H]
\centering
\caption{Distribution of Relevant Records by Data Source}
\label{tab:source_distribution}
\begin{tabular}{lrr}
\toprule
\textbf{Source} & \textbf{n} & \textbf{\%} \\
\midrule
YouTube Comments & 4,457 & 52.8 \\
Google Play Store & 1,460 & 17.3 \\
Apple App Store & 1,303 & 15.5 \\
Reddit & 869 & 10.3 \\
Trustpilot & 346 & 4.1 \\
\midrule
\textbf{Total} & \textbf{8,435} & \textbf{100.0} \\
\bottomrule
\end{tabular}
\end{table}

\subsection{Sentiment Analysis Results}

\subsubsection{Overall Sentiment Distribution}

The overall sentiment distribution revealed a polarized discourse with slightly more negative than positive sentiment (Table~\ref{tab:sentiment_overall}, Figure~\ref{fig:sentiment_distribution}).

\begin{table}[H]
\centering
\caption{Overall Sentiment Distribution}
\label{tab:sentiment_overall}
\begin{tabular}{lrr}
\toprule
\textbf{Sentiment} & \textbf{n} & \textbf{\%} \\
\midrule
Negative & 3,796 & 45.0 \\
Positive & 3,567 & 42.3 \\
Neutral & 1,072 & 12.7 \\
\midrule
\textbf{Total} & \textbf{8,435} & \textbf{100.0} \\
\bottomrule
\end{tabular}
\end{table}

The mean VADER compound score was $M = -0.009$ ($SD = 0.618$), indicating slight overall negativity with high variance, characteristic of polarized opinions.

\begin{figure}[H]
    \centering
    \includegraphics[width=0.75\textwidth]{sentiment_distribution.pdf}
    \caption{Overall sentiment distribution showing polarized discourse with 45.0\% negative, 42.3\% positive, and 12.7\% neutral sentiment.}
    \label{fig:sentiment_distribution}
\end{figure}

\subsubsection{Sentiment by Platform}

Platform-specific sentiment analysis revealed significant variations, with Instagram showing the most negative sentiment and TikTok showing slightly positive sentiment (Table~\ref{tab:sentiment_platform}, Figure~\ref{fig:sentiment_platform}).

\begin{table}[H]
\centering
\caption{Sentiment Statistics by Platform}
\label{tab:sentiment_platform}
\begin{tabular}{lrrrrrrr}
\toprule
\textbf{Platform} & \textbf{n} & \textbf{Mean} & \textbf{SD} & \textbf{Min} & \textbf{Max} & \textbf{Q25} & \textbf{Q75} \\
\midrule
TikTok & 4,181 & 0.025 & 0.608 & $-1.00$ & 1.00 & $-0.51$ & 0.57 \\
Social Media General & 1,654 & 0.007 & 0.627 & $-1.00$ & 1.00 & $-0.56$ & 0.57 \\
YouTube & 941 & $-0.048$ & 0.604 & $-1.00$ & 1.00 & $-0.57$ & 0.44 \\
Facebook & 648 & $-0.076$ & 0.624 & $-1.00$ & 1.00 & $-0.62$ & 0.48 \\
Instagram & 1,011 & $-0.096$ & 0.639 & $-1.00$ & 1.00 & $-0.69$ & 0.46 \\
\bottomrule
\end{tabular}
\end{table}

\begin{figure}[H]
    \centering
    \includegraphics[width=0.85\textwidth]{sentiment_by_platform.pdf}
    \caption{Sentiment comparison across platforms. Instagram exhibits the most negative sentiment ($M = -0.096$), while TikTok shows slightly positive sentiment ($M = 0.025$) despite having the highest discussion volume.}
    \label{fig:sentiment_platform}
\end{figure}

\textbf{Key Finding:} Instagram exhibited the most negative sentiment regarding content concerns, while TikTok---despite receiving the most discussion---showed slightly positive overall sentiment.

\subsubsection{Sentiment by Keyword Category}

Analysis of sentiment by keyword category revealed that parental control discussions had the most negative sentiment, while algorithm-related discussions showed more positive sentiment (Table~\ref{tab:sentiment_keywords}).

\begin{table}[H]
\centering
\caption{Sentiment by Keyword Category}
\label{tab:sentiment_keywords}
\small
\begin{tabular}{lrrrrrr}
\toprule
\textbf{Category} & \textbf{With KW (n)} & \textbf{Mean} & \textbf{SD} & \textbf{Without KW Mean} & \textbf{Difference} \\
\midrule
Parental Controls & 399 & $-0.186$ & 0.557 & 0.026 & $-0.212$ \\
Explicit Content & 210 & $-0.143$ & 0.568 & $-0.014$ & $-0.129$ \\
Mental Health & 130 & $-0.011$ & 0.633 & $-0.035$ & 0.024 \\
Age Concerns & 405 & 0.002 & 0.587 & $-0.047$ & 0.049 \\
Algorithm & 457 & 0.088 & 0.612 & $-0.090$ & 0.178 \\
\bottomrule
\end{tabular}
\end{table}

\subsubsection{Sentiment by Inferred Age Group}

Among records where age could be inferred from textual indicators, teens showed the most negative sentiment (Table~\ref{tab:sentiment_age}, Figure~\ref{fig:sentiment_age}).

\begin{table}[H]
\centering
\caption{Sentiment by Inferred Age Category}
\label{tab:sentiment_age}
\begin{tabular}{lrrrr}
\toprule
\textbf{Age Category} & \textbf{n} & \textbf{Mean} & \textbf{SD} & \textbf{Q25--Q75} \\
\midrule
Teen (13--17) & 50 & $-0.186$ & 0.582 & $-0.67$ to 0.35 \\
Adult (25--34) & 119 & 0.001 & 0.585 & $-0.47$ to 0.50 \\
Parent (35+) & 22 & 0.047 & 0.661 & $-0.62$ to 0.72 \\
Young Adult (18--24) & 23 & 0.232 & 0.560 & $-0.32$ to 0.74 \\
Unknown & 1,205 & $-0.037$ & 0.592 & $-0.55$ to 0.48 \\
\bottomrule
\end{tabular}
\end{table}

\begin{figure}[H]
    \centering
    \includegraphics[width=0.8\textwidth]{sentiment_by_age.pdf}
    \caption{Sentiment distribution by inferred age category. Teens show the most negative sentiment ($M = -0.186$), while young adults show the most positive ($M = 0.232$).}
    \label{fig:sentiment_age}
\end{figure}

\subsection{Topic Modeling Results}

LDA analysis identified five dominant themes in the discourse. Model selection was based on coherence score optimization (Table~\ref{tab:coherence}).

\begin{table}[H]
\centering
\caption{Topic Model Coherence Scores}
\label{tab:coherence}
\begin{tabular}{cc}
\toprule
\textbf{Number of Topics} & \textbf{Coherence Score ($C_v$)} \\
\midrule
5 & \textbf{0.349} \\
7 & 0.335 \\
10 & 0.304 \\
12 & 0.310 \\
15 & 0.286 \\
\bottomrule
\end{tabular}
\end{table}

The five identified topics and their top keywords are presented in Table~\ref{tab:topics}.

\begin{table}[H]
\centering
\caption{LDA Topic Model Results: Five Identified Themes}
\label{tab:topics}
\small
\begin{tabular}{clL{9cm}}
\toprule
\textbf{Topic} & \textbf{Label} & \textbf{Top Keywords (with weights)} \\
\midrule
1 & Child Safety \& Parental Concerns & people (0.022), kid (0.020), child (0.016), social (0.015), medium (0.015), parent (0.008), addiction (0.006), mental (0.006) \\
\addlinespace
2 & Algorithm \& Content Issues & content (0.027), post (0.013), algorithm (0.012), view (0.012), watch (0.010), show (0.009), platform (0.008), block (0.007) \\
\addlinespace
3 & Personal Usage Experiences & year (0.015), said (0.011), day (0.011), time (0.011), friend (0.007), feel (0.007), told (0.006), first (0.007) \\
\addlinespace
4 & Account Moderation \& Banning & account (0.038), reason (0.023), banned (0.016), ban (0.009), community (0.009), support (0.008), channel (0.007) \\
\addlinespace
5 & Mental Health \& Screen Time & time (0.013), social (0.008), brain (0.006), hour (0.006), dopamine (0.004), screen (0.005), attention (0.004) \\
\bottomrule
\end{tabular}
\end{table}

\subsection{Keyword Category Prevalence}

Among the 8,435 relevant records, keyword category occurrence was as follows (Table~\ref{tab:keyword_prevalence}, Figure~\ref{fig:keyword_analysis}):

\begin{table}[H]
\centering
\caption{Keyword Category Prevalence in Relevant Records}
\label{tab:keyword_prevalence}
\begin{tabular}{lrr}
\toprule
\textbf{Category} & \textbf{n} & \textbf{\% of Relevant} \\
\midrule
Age Concerns & 1,791 & 21.2 \\
Parental Controls & 1,554 & 18.4 \\
Algorithm & 1,033 & 12.2 \\
Explicit Content & 610 & 7.2 \\
Mental Health & 353 & 4.2 \\
\bottomrule
\end{tabular}
\end{table}

\begin{figure}[H]
    \centering
    \includegraphics[width=0.8\textwidth]{keyword_category_analysis.pdf}
    \caption{Keyword category prevalence showing age concerns (21.2\%) as the most discussed topic, followed by parental controls (18.4\%) and algorithm-related concerns (12.2\%).}
    \label{fig:keyword_analysis}
\end{figure}

\textbf{Note:} Categories are not mutually exclusive; records may contain keywords from multiple categories.

\subsection{Word Cloud Visualizations}

Figure~\ref{fig:wordclouds} presents word clouds for overall content and platform-specific discussions, visualizing the most frequent terms in the discourse.

\begin{figure}[H]
    \centering
    \begin{subfigure}[b]{0.45\textwidth}
        \includegraphics[width=\textwidth]{wordcloud_all.pdf}
        \caption{All platforms combined}
    \end{subfigure}
    \hfill
    \begin{subfigure}[b]{0.45\textwidth}
        \includegraphics[width=\textwidth]{wordcloud_tiktok.pdf}
        \caption{TikTok}
    \end{subfigure}

    \vspace{0.5cm}

    \begin{subfigure}[b]{0.45\textwidth}
        \includegraphics[width=\textwidth]{wordcloud_instagram.pdf}
        \caption{Instagram}
    \end{subfigure}
    \hfill
    \begin{subfigure}[b]{0.45\textwidth}
        \includegraphics[width=\textwidth]{wordcloud_youtube.pdf}
        \caption{YouTube}
    \end{subfigure}

    \caption{Word clouds showing frequent terms in content-concern discussions. Terms like ``kid,'' ``child,'' ``content,'' and ``algorithm'' appear prominently across platforms.}
    \label{fig:wordclouds}
\end{figure}

\begin{figure}[H]
    \centering
    \begin{subfigure}[b]{0.45\textwidth}
        \includegraphics[width=\textwidth]{wordcloud_positive.pdf}
        \caption{Positive sentiment records}
    \end{subfigure}
    \hfill
    \begin{subfigure}[b]{0.45\textwidth}
        \includegraphics[width=\textwidth]{wordcloud_negative.pdf}
        \caption{Negative sentiment records}
    \end{subfigure}
    \caption{Word clouds comparing positive and negative sentiment records, revealing distinct vocabulary patterns.}
    \label{fig:wordclouds_sentiment}
\end{figure}

\subsection{Correlation Analysis}

Figure~\ref{fig:correlation} shows the correlation heatmap between numerical variables in the dataset.

\begin{figure}[H]
    \centering
    \includegraphics[width=0.75\textwidth]{correlation_heatmap.pdf}
    \caption{Correlation heatmap showing relationships between engagement metrics, sentiment scores, and keyword category indicators.}
    \label{fig:correlation}
\end{figure}

% =============================================================================
% 5. DISCUSSION
% =============================================================================
\section{Discussion}

\subsection{Principal Findings}

This study examined public discourse surrounding content concerns on short-form video platforms through analysis of 8,435 relevant user-generated records from 221,764 collected. Several key findings emerged:

\subsubsection{Finding 1: Polarized but Slightly Negative Discourse}

The overall sentiment distribution (45\% negative, 42.3\% positive) indicates a deeply divided public opinion regarding platform content concerns. The slight negative skew ($M = -0.009$) suggests that users expressing concerns about content issues tend to frame their experiences negatively, though substantial positive discourse exists, potentially representing defensive reactions or satisfaction with platform features. This polarization aligns with broader societal debates about social media's role \citep{haidt2023social}.

\subsubsection{Finding 2: Platform-Specific Sentiment Patterns}

The finding that Instagram exhibited the most negative sentiment ($-0.096$) while TikTok showed slightly positive sentiment ($+0.025$) is noteworthy. This pattern may reflect:

\begin{itemize}[leftmargin=*, itemsep=0.3em]
    \item TikTok's more developed content moderation tools and parental controls \citep{tiktok2023safety}
    \item Instagram's algorithmic integration of Reels into established user bases with different expectations
    \item Demographic differences in platform user populations expressing concerns
    \item Different expectations between native short-form platforms and those that added the feature later
\end{itemize}

\subsubsection{Finding 3: Child Safety as Dominant Concern}

The prevalence of age-related concerns (21.2\% of relevant content) and the prominence of child-related terms in topic modeling confirm that parental and child safety concerns dominate public discourse. This aligns with regulatory attention and media coverage focusing on minor protection \citep{ofcom2023children}. The emergence of a distinct ``Child Safety \& Parental Concerns'' topic (Topic 1) underscores the centrality of this issue.

\subsubsection{Finding 4: Algorithm Awareness and Criticism}

The emergence of a distinct algorithm-focused topic (Topic 2) and the 12.2\% prevalence of algorithm-related keywords indicates significant public awareness of and concern about recommendation systems. Users explicitly discuss how ``For You'' pages expose them to unwanted content, supporting previous research on algorithmic content curation concerns \citep{pariser2011filter, ribeiro2020auditing}.

\subsubsection{Finding 5: Teen Vulnerability}

Among records with inferable age indicators, teens showed the most negative sentiment ($M = -0.186$), while young adults showed the most positive ($M = 0.232$). This pattern supports developmental theories suggesting differential vulnerability across age groups and aligns with concerns about adolescent mental health impacts \citep{twenge2018increases}.

\subsection{Implications for Nepal and Developing Economies}

Our findings have particular relevance for Nepal and similar developing digital economies:

\textbf{1. Digital Literacy Gap:} The prominence of algorithm-related concerns (12.2\%) suggests that even global users struggle to understand content curation. In Nepal, where digital literacy education is limited, this knowledge gap may be even more pronounced, leaving users---especially youth---less equipped to critically evaluate their content exposure.

\textbf{2. Parental Mediation Challenges:} The dominance of child safety concerns (21.2\%) highlights the critical role of parental oversight. However, in developing economies:
\begin{itemize}[itemsep=0.2em]
    \item Parents often have lower digital literacy than children
    \item Extended family structures mean multiple caregivers with varying tech awareness
    \item Economic pressures limit time for active digital mediation
    \item Parental control features may be underutilized due to language barriers
\end{itemize}

\textbf{3. Content Moderation Gaps:} Our Topic 4 (Account Moderation \& Banning) reveals user frustration with moderation systems. For Nepali users:
\begin{itemize}[itemsep=0.2em]
    \item Nepali-language content may be inadequately moderated
    \item Cultural context for ``inappropriate'' content differs from Western norms
    \item Appeal processes may be inaccessible due to language barriers
    \item Over-moderation of legitimate local content may occur
\end{itemize}

\textbf{4. Mental Health Infrastructure:} The mental health discourse (Topic 5) is concerning given Nepal's limited mental health resources:
\begin{itemize}[itemsep=0.2em]
    \item Fewer than 0.5 psychiatrists per 100,000 population
    \item Mental health stigma remains significant
    \item School counseling services are rare
    \item If social media exacerbates mental health issues, support systems are inadequate
\end{itemize}

\textbf{5. Policy Implications:} Nepal's regulatory framework for digital platforms is nascent. Our findings suggest priorities:
\begin{itemize}[itemsep=0.2em]
    \item Mandatory Nepali-language content moderation
    \item Age verification requirements for platform access
    \item Digital literacy curriculum in schools
    \item Parental education programs
    \item Local research funding to understand Nepal-specific effects
\end{itemize}

\subsection{Comparison with Prior Research}

Table~\ref{tab:comparison_findings} compares our key findings with prior research:

\begin{table}[H]
\centering
\caption{Comparison of Current Study Findings with Prior Research}
\label{tab:comparison_findings}
\small
\begin{tabular}{L{3cm}L{5cm}L{5cm}}
\toprule
\textbf{Aspect} & \textbf{Prior Research} & \textbf{This Study} \\
\midrule
Overall effect & Small negative effect, $r = -0.08$ \citep{nguyen2025feeds} & 45\% negative sentiment; $M = -0.009$ \\
\addlinespace
Content specificity & Gap identified; few studies & Explicit content 7.2\% of concerns \\
\addlinespace
Age differences & Adolescents more vulnerable \citep{twenge2018increases} & Teens most negative ($M = -0.186$) \\
\addlinespace
Platform comparison & Limited multi-platform studies & Instagram most negative; TikTok positive \\
\addlinespace
Algorithm concerns & Growing awareness \citep{ribeiro2020auditing} & 12.2\% discuss algorithms explicitly \\
\bottomrule
\end{tabular}
\end{table}

\subsection{Theoretical Implications}

These findings contribute to several theoretical frameworks:

\textbf{Uses and Gratifications Theory:} The polarized sentiment suggests that platforms simultaneously satisfy and frustrate user needs, with content concerns representing gratification barriers \citep{katz1973uses}.

\textbf{Cultivation Theory:} The prominence of mental health discourse (Topic 5) and keywords like ``dopamine'' and ``brain rot'' indicates public internalization of media effects narratives \citep{gerbner1998cultivation}.

\textbf{Social Learning Theory:} Concerns about children's exposure reflect parental awareness of modeling effects and content influence on developing minds \citep{bandura1977social}.

\subsection{Practical Implications}

\subsubsection{For Platform Developers}

\begin{itemize}[leftmargin=*, itemsep=0.3em]
    \item Invest in more granular content filtering options
    \item Improve transparency in algorithmic recommendations
    \item Strengthen age verification mechanisms beyond self-declaration
    \item Develop clearer, more accessible parental control interfaces
    \item \textbf{Invest in local-language content moderation for developing markets}
    \item \textbf{Provide parental controls with interfaces in local languages}
\end{itemize}

\subsubsection{For Parents and Educators}

\begin{itemize}[leftmargin=*, itemsep=0.3em]
    \item Engage in active mediation rather than restrictive-only approaches \citep{livingstone2017parental}
    \item Utilize available parental control features
    \item Discuss algorithmic content curation with young users
    \item Be aware of platform-specific differences in content exposure
\end{itemize}

\subsubsection{For Policymakers (Particularly in Nepal)}

\begin{itemize}[leftmargin=*, itemsep=0.3em]
    \item Consider platform-specific regulatory approaches given sentiment variations
    \item Mandate algorithmic transparency requirements
    \item Support digital literacy education initiatives in schools
    \item Require robust age verification for explicit content access
    \item \textbf{Develop Nepal-specific content moderation guidelines}
    \item \textbf{Fund local research on social media effects}
    \item \textbf{Create parental awareness programs}
\end{itemize}

\subsection{Limitations}

Several limitations should be acknowledged:

\begin{enumerate}[leftmargin=*, itemsep=0.3em]
    \item \textbf{Selection Bias:} Users who leave reviews or post about concerns may not represent typical platform users; those with strong opinions are overrepresented
    \item \textbf{Temporal Scope:} Data represents a cross-sectional snapshot; longitudinal patterns are not captured
    \item \textbf{Geographic Limitations:} English-language content dominates the sample, potentially underrepresenting South Asian perspectives despite the research focus
    \item \textbf{Sentiment Analysis Constraints:} VADER may not capture nuanced, sarcastic, or non-English expressions
    \item \textbf{Causal Inference:} This descriptive study cannot establish causal relationships between platform use and outcomes
    \item \textbf{Age Inference:} Inferred age categories based on textual cues may be inaccurate
    \item \textbf{Nepal-Specific Data:} While motivated by Nepal context, the dataset is not Nepal-specific, limiting direct applicability
\end{enumerate}

\subsection{Future Research Directions}

Future research should:

\begin{itemize}[leftmargin=*, itemsep=0.3em]
    \item \textbf{Conduct primary survey research in Nepal} with diverse demographic samples
    \item \textbf{Collect Nepali-language data} from local platforms and social media
    \item Perform longitudinal tracking of content concern evolution
    \item Investigate cultural variations in platform perceptions, especially South Asian contexts
    \item Develop and test intervention studies for parental mediation strategies
    \item Examine platform policy changes and their effects on user discourse
    \item Apply more sophisticated NLP methods (transformer-based sentiment analysis, aspect-based sentiment analysis)
    \item \textbf{Study the effectiveness of digital literacy interventions in Nepali schools}
\end{itemize}

% =============================================================================
% 6. CONCLUSION
% =============================================================================
\section{Conclusion}

This large-scale secondary data analysis of 221,764 user-generated records reveals that public discourse surrounding short-form video platforms exhibits significant concern regarding content exposure, particularly for minors. The polarized sentiment patterns, platform-specific variations, and emergence of distinct thematic concerns highlight the complexity of the ``For You page problem.''

Key conclusions include:

\begin{enumerate}[leftmargin=*, itemsep=0.5em]
    \item \textbf{Concern is widespread but contested:} Nearly equal proportions of negative (45\%) and positive (42.3\%) sentiment indicate ongoing debate about platform impacts

    \item \textbf{Platform matters:} Instagram generates more negative discourse than TikTok despite TikTok's larger discussion volume, suggesting platform-specific user experiences warrant differentiated approaches

    \item \textbf{Children are central to concerns:} Age-related content dominates the discourse (21.2\%), confirming the salience of child protection as a policy priority

    \item \textbf{Algorithm awareness is growing:} Users increasingly recognize and critique recommendation systems' roles in content exposure, with 12.2\% explicitly discussing algorithms

    \item \textbf{Mental health framing persists:} Discourse incorporates neurological and psychological language (``dopamine,'' ``brain rot,'' ``attention span''), indicating public engagement with scientific narratives about social media effects
\end{enumerate}

\textbf{For Nepal and developing digital economies}, these findings are particularly critical. The rapid adoption of short-form video platforms is outpacing the development of digital literacy, regulatory frameworks, and mental health support systems. The concerns documented in global discourse are likely amplified in contexts where:
\begin{itemize}[itemsep=0.2em]
    \item Parents have lower digital literacy than their children
    \item Content moderation in local languages is inadequate
    \item Mental health resources are scarce
    \item Regulatory mechanisms are nascent
\end{itemize}

This study provides empirical grounding for ongoing policy debates, platform design decisions, and parental guidance regarding short-form video consumption. As these platforms continue to evolve, monitoring public discourse offers valuable insights into emerging concerns and user needs. The identification of content-specific effects as a research gap, now partially addressed by this study, should encourage further investigation---particularly in developing economies like Nepal---into how particular types of content affect users across different age groups and cultural contexts.

\textbf{We call on researchers, policymakers, and platforms to prioritize the unique challenges facing developing digital economies, where the potential harms of inadequate content moderation and digital literacy may be most severe.}

% =============================================================================
% ACKNOWLEDGMENTS
% =============================================================================
\section*{Acknowledgments}

The author thanks the open-source community for providing the data collection and analysis tools used in this research. This study was conducted as independent research from Nepal with no external funding, demonstrating that impactful research is possible even with limited resources.

% =============================================================================
% DATA AVAILABILITY
% =============================================================================
\section*{Data Availability Statement}

The anonymized dataset and analysis scripts are available upon request from the corresponding author. Raw data cannot be shared due to platform terms of service, but collection scripts are provided for replication.

% =============================================================================
% CONFLICT OF INTEREST
% =============================================================================
\section*{Conflict of Interest}

The author declares no conflict of interest.

% =============================================================================
% REFERENCES
% =============================================================================
\newpage
\bibliographystyle{plainnat}
\bibliography{references}

% =============================================================================
% APPENDICES
% =============================================================================
\newpage
\appendix

\section{Supplementary Word Clouds by Keyword Category}

\begin{figure}[H]
    \centering
    \begin{subfigure}[b]{0.45\textwidth}
        \includegraphics[width=\textwidth]{wordcloud_explicit_content.pdf}
        \caption{Explicit content discussions}
    \end{subfigure}
    \hfill
    \begin{subfigure}[b]{0.45\textwidth}
        \includegraphics[width=\textwidth]{wordcloud_age_concerns.pdf}
        \caption{Age-related concerns}
    \end{subfigure}

    \vspace{0.5cm}

    \begin{subfigure}[b]{0.45\textwidth}
        \includegraphics[width=\textwidth]{wordcloud_mental_health.pdf}
        \caption{Mental health discussions}
    \end{subfigure}
    \hfill
    \begin{subfigure}[b]{0.45\textwidth}
        \includegraphics[width=\textwidth]{wordcloud_algorithm.pdf}
        \caption{Algorithm-related discussions}
    \end{subfigure}

    \vspace{0.5cm}

    \begin{subfigure}[b]{0.45\textwidth}
        \includegraphics[width=\textwidth]{wordcloud_parental_controls.pdf}
        \caption{Parental control discussions}
    \end{subfigure}

    \caption{Word clouds for each keyword category showing distinct vocabulary patterns.}
    \label{fig:appendix_wordclouds}
\end{figure}

\section{Data Collection and Analysis Code}

All data collection and analysis scripts are available at: [GitHub Repository - to be added upon publication]

Key scripts include:
\begin{itemize}
    \item \texttt{01a\_scrape\_playstore.py} -- Google Play Store review collection
    \item \texttt{01b\_scrape\_appstore.py} -- Apple App Store review collection
    \item \texttt{01c\_scrape\_reddit.py} -- Reddit post collection
    \item \texttt{01h\_combine\_all\_data.py} -- Data combination and standardization
    \item \texttt{03\_sentiment\_analysis.py} -- VADER sentiment analysis
    \item \texttt{04\_topic\_modeling.py} -- LDA topic modeling
    \item \texttt{05\_visualizations.py} -- Figure generation
\end{itemize}

\section{Interactive Visualizations}

An interactive pyLDAvis visualization of the topic model is available at: \\
\texttt{outputs/topics/lda\_visualization.html}

\end{document}
